\section{Relative bundle quotients and certificates of strict progress}
\label{sec:quotient-progress}

Lackenby's equality criterion for quadratic handle complexity gives an
algebraic way to certify that a proposed surface makes strict progress.
The useful object is the part of relative second homology that is
represented inside the interior parallelity bundle. A surface whose class
survives after quotienting out that subspace cannot be one of the
homologically bundle-supported surfaces allowed by the equality case.
We make this observation quantitative and give a complete exact
certificate system for its algebraic premise.

The geometric implication below is a corollary of
\cite[Proposition~5.9]{lackenby2026}. The linear alternative is an
application of finite-dimensional duality. The contribution here is their
explicit combination with a small, independently replayable certificate
contract, including the conditions under which the algebraic result can
be promoted to a statement about the geometric algorithm.

\subsection{The pair map and the geometric hypotheses}

Let $M$ be a compact connected oriented PL $3$-manifold with boundary
pattern $P$, and let $\mathcal H$ be an admissible handle structure in the
sense of \cite[Definition~5.3]{lackenby2026}. Let $\mathcal B$ be its
\emph{actual interior parallelity bundle}, as defined in
\cite[Definition~5.6]{lackenby2026}, with the embedding in $M$ inherited
from those handles. Throughout this section use the conservative pair
\[
 A=\mathcal B\cap\partial M,\qquad
 j:(\mathcal B,A)\longrightarrow(M,\partial M).
\]
This is an inclusion of pairs, so it induces a well-defined map on
relative homology. Define
\begin{equation}
 \label{eq:quotient-bundle-image}
 V_{\mathcal B}
 =\operatorname{im}\!\left(
 H_2(\mathcal B,A;\mathbb Q)
 \xrightarrow{\,j_*\,} H_2(M,\partial M;\mathbb Q)
 \right).
\end{equation}
All chain complexes used to compute this map must come from compatible
cellular or simplicial models of the \emph{embedded} pair. An abstract
isomorphism between homology groups does not specify this image.

\begin{remark}[Horizontal boundary is an additional embedding assertion]
\label{rem:quotient-horizontal}
One can instead use $A_h=\partial_h\mathcal B$ if its actual inclusion in
$\partial M$ has been established. For the interior parallelity handles,
having no $3$-handles is a sufficient condition: the horizontal faces are
disjoint from other $0$-, $1$-, and $2$-handles and hence lie in the
manifold boundary. A general handle structure may have $3$-handles, and
horizontal faces may meet them. Thus a map
$H_2(\mathcal B,\partial_h\mathcal B)\to H_2(M,\partial M)$ must not be
formed solely from an abstract description of the bundle.
The choice $A=\mathcal B\cap\partial M$ always defines the pair map needed
here. When $A_h\subseteq A$, its image contains the image obtained from
$(\mathcal B,A_h)$, so it can make the obstruction weaker but does not
invalidate the sufficient implication. Twisted $I$-bundles are allowed.
\end{remark}

For a $0$-handle $H$, write $I(H)$ for its integer index from
\cite[Definition~5.1]{lackenby2026}. The quadratic complexity is
\[
 c_q(\mathcal H)=\sum_{H\in\mathcal H^0}
                  \bigl(\max\{I(H),0\}\bigr)^2.
\]
The subscript $q$ denotes quadratic complexity. Later the certificate
file also has a field named \texttt{q}; we denote that positive integer
by $\nu$ in the mathematics to distinguish the two uses.

\begin{theorem}[The quotient obstruction forces a positive even raw defect]
\label{thm:quotient-strict-progress}
Let $\mathcal H$, $M$, $P$, $\mathcal B$, and $A$ be as above. Let $S$ be
a compact connected oriented surface properly embedded in $M$ and normal
with respect to $\mathcal H$. Suppose
\[
 [S]\notin V_{\mathcal B}.
\]
Let $\mathcal H^{\mathrm{raw}}$ be the handle structure obtained by
cutting along $S$, before any deletion, handle cancellation, merging,
admissibility repair, or bundle replacement. For each original $0$-handle
$H$, let
$H_{H,1},\ldots,H_{H,r_H}$ be \emph{all} its raw child $0$-handles, and put
$a_{H,i}=I(H_{H,i})$. Then
\begin{equation}
 \label{eq:quotient-exact-raw-defect}
 c_q(\mathcal H)-c_q(\mathcal H^{\mathrm{raw}})
 =2\sum_{H\in\mathcal H^0}\ \sum_{1\leq i<j\leq r_H}a_{H,i}a_{H,j}
 \ \geq\ 2.
\end{equation}
If $\mathcal H^{\mathrm{clean}}$ is subsequently obtained by operations
independently proved not to increase $c_q$, then
\[
 c_q(\mathcal H)-c_q(\mathcal H^{\mathrm{clean}})\geq 2
\]
as well. The exact identity in
\eqref{eq:quotient-exact-raw-defect} concerns the raw cut.
\end{theorem}

\begin{proof}
Since $0\in V_{\mathcal B}$, the hypothesis implies $[S]\neq0$ in
$H_2(M,\partial M;\mathbb Q)$. Orientability of $M$ and $S$ makes $S$
two-sided. If $S$ separated $M$, the closure of one side would give a
relative $3$-chain whose boundary is $S$, up to orientation, together
with a chain in $\partial M$. This would imply $[S]=0$. Consequently
$S$ is nonseparating, and the hypotheses of
\cite[Proposition~5.9]{lackenby2026} hold.

Admissibility implies $I(H)\geq0$. Index conservation under a standard
cut and nonnegativity after a normal cut of an admissible handle
structure give
\[
 \sum_{i=1}^{r_H}a_{H,i}=I(H),\qquad a_{H,i}\geq0
\]
for every $H$; these are
\cite[Lemmas~5.2 and~5.5]{lackenby2026}. Expanding the square gives
\[
 I(H)^2-\sum_i a_{H,i}^2
 =\left(\sum_i a_{H,i}\right)^2-\sum_i a_{H,i}^2
 =2\sum_{i<j}a_{H,i}a_{H,j}.
\]
Summing over all original $0$-handles proves the exact formula and shows
that the raw defect is an even nonnegative integer.

Suppose that the defect were zero. The equality conclusion of
\cite[Proposition~5.9]{lackenby2026} would make $S$ homologous in
$H_2(M,\partial M)$ to a vertical surface contained in $\mathcal B$.
The proof of that proposition arranges the representative to be properly
embedded in $\mathcal B$, disjoint from $\partial_v\mathcal B$, and with
boundary in $\partial M$. In particular its boundary lies in
$A=\mathcal B\cap\partial M$, so it defines a relative class in
$H_2(\mathcal B,A;\mathbb Q)$ whose image is $[S]$.
This contradicts $[S]\notin V_{\mathcal B}$. The even nonnegative defect
is therefore at least two.

For the last assertion, write the total defect as
\[
 \bigl(c_q(\mathcal H)-c_q(\mathcal H^{\mathrm{raw}})\bigr)
 +\bigl(c_q(\mathcal H^{\mathrm{raw}})
                  -c_q(\mathcal H^{\mathrm{clean}})\bigr).
\]
The first summand is at least two and the second is nonnegative by the
separately established property of the cleanup.
\end{proof}

The formula must retain the complete raw child inventory. A parallelity
$0$-handle with certified index zero contributes zero and may be omitted
from a numerical sum after that fact is established. An already reduced
inventory may also have lost positive-index pieces or changed the parent
relation, so the conservation identity cannot simply be applied to it.
There is no asserted parity for an additional cleanup defect.
Irreducibility of $M$ and essentiality of $P$ are not hypotheses of
Proposition~5.9 or of Theorem~\ref{thm:quotient-strict-progress}. They must
still be supplied if a surrounding cleanup or hierarchy routine requires
them. For disconnected $M$, apply the argument in the component
containing $S$ and add the unchanged complexities of the other components.

\begin{corollary}[A bound for a retained monotone sequence]
\label{cor:quotient-retained-count}
Consider a retained forward sequence of handle structures starting at
$\mathcal H_0$. Suppose every intervening operation is
$c_q$-nonincreasing and $N$ of the cuts satisfy
Theorem~\ref{thm:quotient-strict-progress}. Then
\[
 N\leq
 \left\lfloor\frac{c_q(\mathcal H_0)-c_q(\mathcal H_{\mathrm{last}})}2
       \right\rfloor
 \leq \left\lfloor\frac{c_q(\mathcal H_0)}2\right\rfloor .
\]
\end{corollary}

\begin{proof}
Each of the $N$ certified cuts contributes a drop of at least two.
All other contributions are nonnegative, so the drops telescope.
\end{proof}

This corollary counts certified cuts on the retained monotone sequence.
It does not count discarded attempts, restored states, or later
reconstructions after backtracking. It gives additive progress and
therefore does not establish logarithmic hierarchy depth. A proof
controlling the complete search must also charge the work on histories
that are subsequently discarded.

\subsection{A complete finite-chain alternative}

We now isolate an algebraic problem that can be checked without trusting
a proposed basis of bundle cycles. Let the dimensions of $C_i$ and $E_i$
be $c_i$ and $e_i$. The supplied rational matrices are
\[
 d_2:C_2\longrightarrow C_1,\quad
 d_3:C_3\longrightarrow C_2,\quad
 \delta_2:E_2\longrightarrow E_1,\quad
 j_1:E_1\longrightarrow C_1,\quad
 j_2:E_2\longrightarrow C_2.
\]
The input also includes $s\in C_2$ and must satisfy
\begin{equation}
 \label{eq:quotient-input-contract}
 d_2d_3=0,\qquad d_2j_2=j_1\delta_2,\qquad d_2s=0.
\end{equation}
These equations suffice to define
\[
 H=\ker d_2/\operatorname{im}d_3,\qquad
 V=\bigl(\operatorname{im}d_3+j_2\ker\delta_2\bigr)/
       \operatorname{im}d_3\ \subseteq H.
\]
If the matrices arise from authenticated relative chain complexes and
the actual inclusion of pairs, then $H=H_2(M,\partial M;\mathbb Q)$ and
$V=V_{\mathcal B}$. Without that authentication, they define precisely
the displayed algebraic quotient and no asserted manifold.

\begin{theorem}[Primal support or dual obstruction]
\label{thm:quotient-alternative}
For data satisfying \eqref{eq:quotient-input-contract}, exactly one of
the following systems has a rational solution:
\begin{align}
 \text{\rm(P)}\qquad&
 s=d_3a+j_2b,\qquad \delta_2b=0,
 &&a\in C_3,\quad b\in E_2; \label{eq:quotient-primal}\\
 \text{\rm(D)}\qquad&
 yd_3=0,\qquad yj_2=u\delta_2,\qquad ys=1,
 &&y\in C_2^*,\quad u\in E_1^* . \label{eq:quotient-dual}
\end{align}
System {\rm(P)} holds exactly when $[s]\in V$; system {\rm(D)} holds
exactly when $[s]\notin V$.
\end{theorem}

\begin{proof}
By the definition of $V$, membership is equivalent to the existence of
$b\in\ker\delta_2$ such that $s-j_2b\in\operatorname{im}d_3$.
This is precisely (P). If (P) and (D) both held, then
\[
 1=ys=yd_3a+yj_2b=u\delta_2b=0,
\]
which is impossible.

To prove existence of (D) when (P) fails, take a matrix $Z$ whose columns
are a basis of $\ker\delta_2$ and put $W=[\,d_3\mid j_2Z\,]$.
Then $s\notin\operatorname{col}W$. Define a linear functional on
$\operatorname{span}(\operatorname{col}W,s)$ to vanish on
$\operatorname{col}W$ and to take value one on $s$, and extend it to
$C_2$. Call its row vector $y$. It satisfies $yd_3=0$, $ys=1$, and
$yj_2Z=0$.

For completeness, a row $v$ annihilates $\ker\delta_2$ if and only if
it is in the row space of $\delta_2$. One direction follows by
multiplication. For the converse, the rule
$\delta_2b\mapsto vb$ is well-defined on $\operatorname{im}\delta_2$
when $v$ vanishes on its kernel; extending this functional to $E_1$
gives a row $u$ with $v=u\delta_2$. Apply this fact to $v=yj_2$.
It gives $yj_2=u\delta_2$, completing (D).
\end{proof}

The matrix $Z$ occurs only in this proof. The certificate does not accept
a purported cycle basis as trusted input. A basis with missing vectors
could make $yj_2Z=0$ hold while $y$ fails to annihilate the images of
other bundle cycles. The identity $yj_2=u\delta_2$ checks annihilation
of \emph{every} bundle cycle directly.

Only the displayed chain fragment is needed for the algebraic
calculation. If a full inclusion $j_*:E_*\to C_*$ is supplied, the
boundaries in $E_2$ already map to boundaries in $C_2$ through
$j_2\delta_3=d_3j_3$; enumerating $E_3$ is unnecessary for computing
the image from all cycles of $E_2$. This observation does not dispense
with authenticating the full geometric map when that interpretation is
claimed.

\begin{corollary}[Integer certificate form]
\label{cor:quotient-integer-alternative}
For integral input matrices and integral $s$, exactly one of the
following types of integer certificate exists:
\begin{align}
 \text{\rm(P$_{\mathbb Z}$)}\quad&
 \nu s=d_3a+j_2b,\quad \delta_2b=0,
 &&\nu\in\mathbb Z_{>0}; \label{eq:quotient-integer-primal}\\
 \text{\rm(D$_{\mathbb Z}$)}\quad&
 Yd_3=0,\quad Yj_2=U\delta_2,\quad Ys=\nu,
 &&\nu\in\mathbb Z_{>0}. \label{eq:quotient-integer-dual}
\end{align}
Here $a,b,Y,U$ have integral entries and the same dimensions as above.
\end{corollary}

\begin{proof}
Clear common denominators in a rational certificate, including the
coefficient one on the normalized right-hand side. Conversely divide an
integer certificate by $\nu$. The resulting rational alternatives are
exclusive by Theorem~\ref{thm:quotient-alternative}. Equivalently, if
both integer certificates existed with multipliers $\nu_P,\nu_D>0$,
pairing them would give $\nu_P\nu_D=0$.
\end{proof}

The positive multiplier is essential because the question is over
$\mathbb Q$. For example, take $C_2=E_2=\mathbb Q$, let $j_2$ be
multiplication by $2$, set $s=1$, and let all other relevant spaces be
zero. The class is supported over $\mathbb Q$, with integer certificate
$\nu=2$, $b=1$. Requiring an integer certificate with $\nu=1$ would
confuse rational image membership with integral divisibility.

The auxiliary row $U$ is also substantive. Take
$C_2=E_1=E_2=\mathbb Q$, $\delta_2=[3]$, $j_2=[2]$, $s=1$, and
$C_1=C_3=0$. Since $\ker\delta_2=0$, the class is outside the image.
A certificate is
\[
 Y=[3],\qquad U=[2],\qquad \nu=3.
\]
Its restriction to $E_2$ is a coboundary rather than the zero row.

\subsection{Witness size and exact production}

The matrix system used by the reference producer is
\begin{equation}
 \label{eq:quotient-block-system}
 \mathsf A x=\mathsf f,\qquad
 \mathsf A=
 \begin{pmatrix}d_3&j_2\\0&\delta_2\end{pmatrix},\qquad
 \mathsf f=\begin{pmatrix}s\\0\end{pmatrix},\qquad
 x=\begin{pmatrix}a\\b\end{pmatrix}.
\end{equation}
Let $R=c_2+e_1$, $K=c_3+e_2$, and $r=\operatorname{rank}\mathsf A$.
Suppose the entries of $\mathsf A$ and $\mathsf f$ are integral and
have absolute value at most $M\geq1$. Define
\[
 D_0=1,\qquad D_k=(\sqrt{k}\,M)^k\quad(k\geq1).
\]

\begin{theorem}[Polynomial-size integer witnesses]
\label{thm:quotient-minor-bound}
If \eqref{eq:quotient-block-system} is consistent, a certificate of type
{\rm(P$_{\mathbb Z}$)} exists with
\[
 1\leq\nu\leq D_r,\qquad
 |a_i|,|b_i|\leq D_r.
\]
If the system is inconsistent, a certificate of type
{\rm(D$_{\mathbb Z}$)} exists with
\[
 |Y_i|,|U_i|\leq D_r,\qquad
 1\leq\nu\leq D_{r+1}.
\]
Consequently every entry of either certificate can be encoded with
$O((r+1)(\log(r+2)+\log M))$ bits, and the total certificate length is
polynomial in the explicit chain ranks and input-entry bit lengths.
\end{theorem}

\begin{proof}
In the consistent case, choose $r$ independent columns of $\mathsf A$
and $r$ rows on which they form a nonsingular matrix $T$. A basic
solution sets all other variables to zero. Cramer's rule expresses
each remaining variable as the determinant of $T$ with one column
replaced by the corresponding entries of $\mathsf f$, divided by
$\det T$. Because the chosen columns span the column space and the
system is consistent, this solution satisfies every row, not only
the selected rows. Put $\nu=|\det T|$ and choose the numerator signs
accordingly. Hadamard's inequality bounds every determinant just
mentioned by $D_r$. If $r=0$, consistency forces $\mathsf f=0$, and
one can take $\nu=1$ and all coefficients zero.

In the inconsistent case,
$\operatorname{rank}[\,\mathsf A\mid\mathsf f\,]=r+1$.
Choose $r$ independent columns of $\mathsf A$ and $r+1$ rows on which
those columns together with $\mathsf f$ form a nonsingular
$(r+1)\times(r+1)$ matrix $T^+$. The cofactors of the last column
define an integer row $\lambda$, supported on the selected rows.
Laplace expansion gives
\[
 \lambda\mathsf f=\det T^+\neq0,
 \qquad \lambda\mathsf A=0.
\]
For the second identity, the cofactors annihilate each selected
column by the determinant with two equal columns; they then
annihilate all columns because the selected ones span
$\operatorname{col}\mathsf A$.
Reverse the sign of $\lambda$ if necessary and put
$\nu=\lambda\mathsf f>0$. Splitting the row according to the two
row blocks as $\lambda=(Y,-U)$ gives
\[
 Yd_3=0,\qquad Yj_2-U\delta_2=0,\qquad Ys=\nu.
\]
Each cofactor is an $r\times r$ determinant, so its absolute value is
at most $D_r$; the final determinant has absolute value at most
$D_{r+1}$. This also covers $r=0$, using the determinant of the empty
matrix equal to one. Taking binary lengths proves the last assertion.
\end{proof}

These bounds are existence bounds obtained from selected minors.
They do not require total unimodularity, saturation of an integral
lattice, or a unimodular basis change. Rational input can be reduced
to integral input by clearing denominators; the resulting entry
lengths remain polynomial in the total rational input length.

\begin{proposition}[The row-provenance producer is complete]
\label{prop:quotient-producer}
Over exact rational arithmetic, Gaussian elimination applied to
\eqref{eq:quotient-block-system}, while retaining every elementary row
operation on an identity matrix, produces one of the integer
certificates in Corollary~\ref{cor:quotient-integer-alternative}.
It uses polynomially many arithmetic operations with polynomially
bounded operand lengths.
\end{proposition}

\begin{proof}
Start with the augmented matrix $[\mathsf A\mid\mathsf f]$ and a
provenance matrix $T=I_R$. At every row swap, row rescaling, or row
addition, perform the identical operation on $T$. The maintained
invariant is
\[
 [\,\widetilde{\mathsf A}\mid\widetilde{\mathsf f}\,]
       =T[\,\mathsf A\mid\mathsf f\,].
\]
The producer selects pivots only in coefficient columns and eliminates
each pivot from all other rows.

If a row has zero coefficients and nonzero right-hand side $\beta$,
let $t$ be its row of $T$ and set $\lambda=t/\beta$.
Then $\lambda\mathsf A=0$ and $\lambda\mathsf f=1$.
Writing $\lambda=(y,-u)$ gives
$(yd_3,yj_2-u\delta_2)=(0,0)$ and $ys=1$.
Clearing the denominators of $y$, $u$, and the extra entry $1$
produces integer $(Y,U,\nu)$ with $\nu>0$.
The minus sign in the lower block is necessary and is the sign
used by the reference implementation.

If no inconsistent row occurs, assign zero to every nonpivot
variable and the reduced right-hand-side entry to each pivot
variable. Reduced echelon form then gives a solution $x=(a,b)$.
Clear the denominators of its entries and the extra entry $1$.
This gives $(a,b,\nu)$ satisfying
\eqref{eq:quotient-integer-primal}, with $\nu>0$.
Dividing all certificate entries by their common positive greatest
common divisor preserves the equations. These arguments also apply
when any of the spaces has dimension zero. If $R=0$, the zero vector
of length $K$ with $\nu=1$ is a support certificate.

There are polynomially many row operations on matrices of dimensions
bounded by $R$ and $K+R+1$. With reduced rational arithmetic, the
entries arising in elimination, including the provenance entries,
are ratios of minors of matrices formed from the input and an
identity block. Hadamard bounds their numerators and denominators
by a polynomial number of bits. Products, sums, greatest common
divisors, and common-denominator clearing consequently have
polynomial bit complexity. This proves the claim for the exact
algorithm implemented using \texttt{Fraction}.
\end{proof}

The size convention in the last two statements counts the declared
chain ranks, in addition to matrix-entry bit lengths. It agrees with
a cellular input that explicitly lists its generators. It is not a
claim of polynomial complexity in bare JSON byte length for a
succinct declaration of huge empty matrices: a file can declare
$K$ unconstrained generators using $O(\log K)$ digits, whereas the
current producer writes a vector of length $K$.

\subsection{Independent replay and its exact scope}

The implementation in \texttt{research/quotient} separates a rational
producer from an integer verifier. The input schema records the five
dimensions, all five matrices, and $s$ explicitly; the matrix called
$\delta_2$ above is the JSON field \texttt{e2}. Dimensions are retained
even for empty matrices. The verifier rejects booleans, floating-point
values, strings used as integers, inconsistent shapes, unrecognized
fields, and failures of \eqref{eq:quotient-input-contract}.

\begin{proposition}[Soundness of independent integer replay]
\label{prop:quotient-replay}
For input satisfying \eqref{eq:quotient-input-contract}, the following
replay rules are sound:
\begin{enumerate}
 \item accept a support certificate only after checking strict
 integer types, dimensions, $\nu>0$, and all identities in
 \eqref{eq:quotient-integer-primal};
 \item accept an outside certificate only after checking strict
 integer types, dimensions, $\nu>0$, and all identities in
 \eqref{eq:quotient-integer-dual}.
\end{enumerate}
The first verdict proves $[s]\in V$. The second proves $[s]\notin V$.
Replay uses only exact integer products, sums, and equality checks,
and need not invoke the producer, Gaussian elimination, or a homology
basis computation.
\end{proposition}

\begin{proof}
Dividing the first certificate by $\nu$ gives (P), and hence
membership. For the second, suppose $[s]\in V$. Then
$s=d_3a+j_2b$ for rational $a,b$ with $\delta_2b=0$.
The replayed identities would imply
\[
 \nu=Ys=Yd_3a+Yj_2b=U\delta_2b=0,
\]
contradicting $\nu>0$. The operations stated in the proposition
are exactly those needed to evaluate the identities.
\end{proof}

The verdict strings are \texttt{INSIDE\_BUNDLE\_IMAGE} and
\texttt{OUTSIDE\_BUNDLE\_IMAGE}. For arbitrary input matrices these
refer to the algebraic subspace $V$ defined above. In particular,
an inside verdict does not assert that a geometric cut is neutral:
Lackenby's implication is that equality of complexity forces bundle
support, not that every bundle-supported surface realizes equality.
An outside verdict becomes a strict geometric progress certificate
only after all hypotheses of
Theorem~\ref{thm:quotient-strict-progress} have been authenticated.

The certificate also records a SHA-256 digest of a canonical JSON
serialization of the input. This is useful provenance metadata and
rejects ordinary stale-input mistakes. Algebraic soundness does not
rest on a collision-resistance assumption: the identities are
evaluated against the supplied input itself.

The exact algorithm and the default command-line integer codec must
also be distinguished. This artifact uses JSON decimal integers.
CPython's default limit on conversion to or from very long decimal
strings can reject a mathematically valid large input or output.
The documented launch setting
\texttt{PYTHONINTMAXSTRDIGITS=0} is required for trusted research
instances exceeding that limit. A subprocess test exercises a
$16384$-bit input coefficient, the resulting $16384$-bit support
multiplier, and separate producer and verifier commands under this
setting. The quotient wire format does not use the hexadecimal
codec of the separate raw-cut artifact.

\subsection{The dual class and component selection}

The outside row is not a surface-producing cohomology class of
degree one. If the chain complexes come from the indicated pairs,
$yd_3=0$ says that $y$ is a relative degree-two cocycle, and
$yj_2=u\delta_2$ says that its restriction to the bundle pair is a
coboundary. Thus its class belongs to
\begin{equation}
 \label{eq:quotient-cohomological-annihilator}
 \ker\!\left(
 H^2(M,\partial M;\mathbb Q)
 \longrightarrow H^2(\mathcal B,A;\mathbb Q)
 \right).
\end{equation}
Over $\mathbb Q$, universal coefficients and Poincare--Lefschetz
duality give the identifications
\[
 H^2(M,\partial M;\mathbb Q)
 \cong\operatorname{Hom}_{\mathbb Q}
       \bigl(H_2(M,\partial M;\mathbb Q),\mathbb Q\bigr)
 \cong H_1(M;\mathbb Q).
\]
Consequently the separating functional is represented by an
\emph{absolute one-cycle} $z$, and its value on $[S]$ is the
algebraic intersection of $z$ with $S$.

By contrast, the \emph{surface class} $\sigma=[S]$ is the
Poincare--Lefschetz dual of a class
$\alpha\in H^1(M;\mathbb Q)$.
With compatible sign conventions the relevant evaluation is
$\alpha(z)=y(\sigma)$. It is $\alpha$, rather than $y$, that a
cocycle-to-normal-surface construction should realize. There is
no canonical identification sending the separating row itself
to the required class of $H^1(M;\mathbb Q)$.

For example, in a compatible triangulation with relative
fundamental $3$-cycle $F$, fix a cap-product convention and send
an absolute simplicial $1$-cocycle $c$ to the relative $2$-cycle
$s=c\cap F$. Its class is the dual of $[c]$. This gives a
chain-level route from an authenticated cocycle to the input
of the quotient verifier without expanding an exponentially
large normal surface. Authenticating $F$, the cap product, and
the signed intersection convention remains necessary.

\begin{lemma}[Reuse of a witness on a component]
\label{lem:quotient-component}
Suppose integer rows $Y,U$ satisfy
$Yd_3=0$ and $Yj_2=U\delta_2$.
Let $s_1,\ldots,s_m$ be integral relative cycles, and suppose
\[
 Y\left(\sum_{i=1}^m s_i\right)>0.
\]
Then some $i$ has $Ys_i>0$. For any such $i$, the same rows $Y,U$,
with positive multiplier $\nu_i=Ys_i$, are a valid outside
certificate for $s_i$.
\end{lemma}

\begin{proof}
The displayed positive number is the sum of the integers $Ys_i$,
so at least one summand is positive. The two annihilation identities
do not depend on the candidate cycle. Substituting $s_i$ and
$\nu_i=Ys_i>0$ therefore gives all of
\eqref{eq:quotient-integer-dual}.
\end{proof}

Apply the lemma to the oriented components of a properly embedded
normal representative. It shows that a future compressed component
extractor need only find one component with positive evaluation
under a fixed signed scalar functional; computing a complete
homology vector for every component is unnecessary for this
selection. A normal component with nonzero evaluation has class
outside $V_{\mathcal B}$ and is therefore nonseparating.
A serialized certificate for the selected component must also bind to that component's input digest and use its new positive multiplier; only the annihilating rows are reused. The statement requires the component cycles and their orientations
to be authenticated. Unoriented normal coordinates or unsigned
component weights do not establish the required evaluations.

\subsection{What a polynomial geometric progress step would require}

Suppose the quotient
$H_2(M,\partial M;\mathbb Q)/V_{\mathcal B}$ is nonzero.
At the level of existence, choose a class $\sigma$ surviving in
this quotient, multiply it by a positive integer so that it is
integral, and use its dual class in $H^1(M;\mathbb Z)$.
The cocycle construction in
\cite[Theorem~6.2 and its proof]{lackenby2026} gives a route to a
properly embedded cooriented normal representative. On a
triangulation, one normalizes a cocycle to vanish on a maximal
tree, places signed intersection points on edges, and joins
them by cooriented normal arcs and disks. The resulting total
relative class is the prescribed dual class.
For patterned handle structures, the cell construction and
local normalization require the stated local type bounds;
the theorem additionally assumes no $3$-handles when
$\partial M\neq\varnothing$. This hypothesis must be retained
when invoking that version of the construction.

Class-preserving normalization may split the representative.
Keeping all components preserves its total relative class;
discarding only components with zero relative class is also
safe. At least one component survives the quotient. More
precisely, first obtain an outside witness for the total class
and use Lemma~\ref{lem:quotient-component} to select a component
with positive evaluation. Keeping an arbitrary component is
not justified: it may lie in $V_{\mathcal B}$ even though the
total class does not.

\begin{proposition}[A conditional geometric implementation interface]
\label{prop:quotient-interface}
Assume the following operations are available with polynomial
running time and polynomial output bit budgets in the maintained
geometric representation:
\begin{enumerate}
 \item a compiler and verifier for the actual admissible
 handle structure, its interior parallelity bundle, the
 embedded pair $(\mathcal B,A)\subseteq(M,\partial M)$,
 and the resulting chain maps;
 \item a rational cohomology basis and an explicit
 chain-level Poincare--Lefschetz duality map;
 \item construction and normalization of a prescribed
 integral cohomology class as a compressed cooriented normal
 surface, preserving the total relative homology class
 and satisfying the relevant handle hypotheses;
 \item extraction of a connected oriented component with
 positive evaluation under a supplied signed scalar
 functional whenever the total evaluation is positive,
 together with an authenticated relative cycle for that
 component.
\end{enumerate}
If $H_2(M,\partial M;\mathbb Q)/V_{\mathcal B}$ is nonzero,
these interfaces and the quotient producer and verifier give
a polynomial-time construction of a connected oriented normal
surface certified to have raw $c_q$-defect at least two.
\end{proposition}

\begin{proof}
The duals of a rational basis of $H^1(M;\mathbb Q)$ span
$H_2(M,\partial M;\mathbb Q)$, so at least one lies outside
$V_{\mathcal B}$. Test the finitely many basis vectors by the
exact alternative and choose an outside vector. Clear
denominators with a positive multiplier, retaining the
correspondingly rescaled positive evaluation.
The third interface realizes the resulting integral class
and preserves its total value under the dual witness.
The fourth interface and Lemma~\ref{lem:quotient-component}
select a connected normal component with positive evaluation.
The independent verifier accepts the same annihilating rows
with this new positive multiplier. The authenticated geometric
data then permit Theorem~\ref{thm:quotient-strict-progress}.
The number of basis vectors, the linear algebra, the certificate
sizes, and the four assumed operations all have polynomial
budgets, proving the conditional running-time statement.
\end{proof}

The present artifact implements the finite-chain producer and
replay checker. It does not implement the source-to-chain
compiler, the prescribed-class geometric normalizer, or the
signed compressed component extractor in this proposition.
Weighted orbit methods suggest a route to the last interface,
but a method for unoriented component counting alone does not
supply its authentication or signed evaluations.
The quotient can also vanish; the proposition has no premise
in that case. Resolving such states still requires the
bundle-supported reductions and the other hierarchy mechanisms.

